
\documentclass{article}
\usepackage{colortbl}
\usepackage{makecell}
\usepackage{multirow}
\usepackage{supertabular}

\begin{document}

\newcounter{utterance}

\centering \large Interaction Transcript for game `cladder', experiment `full\_v1.5\_default', episode 1743 with selfrefine{-}llama{-}3.1{-}8b{-}instruct.
\vspace{24pt}

{ \footnotesize  \setcounter{utterance}{1}
\setlength{\tabcolsep}{0pt}
\begin{supertabular}{c@{$\;$}|p{.15\linewidth}@{}p{.15\linewidth}p{.15\linewidth}p{.15\linewidth}p{.15\linewidth}p{.15\linewidth}}
    \# & \multicolumn{2}{c}{Player} && \multicolumn{2}{c}{Game Master} \\
    \hline

    \theutterance \stepcounter{utterance}  
    & & & \multicolumn{4}{p{0.6\linewidth}}{
        \cellcolor[rgb]{0.9,0.9,0.9}{
            \makecell[{{p{\linewidth}}}]{
                \texttt{\tiny{[P1$\langle$GM]}}
                \texttt{You are an expert in causal inference. The following question is not a typical commonsense query, but rather a meticulously designed question created by a professor specializing in causal inference, intended to assess the students' mastery of the course content.} \\
\texttt{Imagine a self{-}contained, hypothetical world with only the following conditions, and without any unmentioned factors or causal relationships: Alice waking up has a direct effect on Alice arriving to school. Traffic has a direct effect on Alice arriving to school. For days when Alice wakes up on time, the probability of arriving to school on time is 46\%. For days when Alice wakes up late, the probability of arriving to school on time is 78\%. For days when Alice wakes up late, would it be more likely to see arriving to school on time if Alice had gotten up on time?} \\
\texttt{Start your answer with 'Yes' or 'No', followed by additional reasoning or evidence to support your explanation.} \\
            }
        }
    }
     \\ \\

    \theutterance \stepcounter{utterance}  
    & & \multicolumn{4}{p{0.6\linewidth}}{
        \cellcolor[rgb]{0.95,0.95,0.95}{
            \makecell[{{p{\linewidth}}}]{
                \texttt{\tiny{[P1$\leftrightarrow$P]}}
                \texttt{{[}DRAFT{]}} \\
\texttt{No.} \\
\\ 
\texttt{To answer this question, we need to apply the concept of counterfactuals and potential outcomes. We are given the following information:} \\
\\ 
\texttt{{-} P(Arriving on time | Waking up on time) = 46\%} \\
\texttt{{-} P(Arriving on time | Waking up late) = 78\%} \\
\\ 
\texttt{We want to find P(Arriving on time | Waking up late, Counterfactual: Waking up on time). This is the probability of arriving on time if Alice had woken up on time, given that she actually woke up late.} \\
\\ 
\texttt{To solve this, we can use the formula for counterfactuals:} \\
\\ 
\texttt{P(Arriving on time | Waking up late, Counterfactual: Waking up on time) = P(Arriving on time | Waking up on time) + P(Arriving on time | Waking up late) {-} P(Arriving on time)} \\
\\ 
\texttt{However, we don't have the value of P(Arriving on time) directly. But we can use the law of total probability to find it:} \\
\\ 
\texttt{P(Arriving on time) = P(Arriving on time | Waking up on time) * P(Waking up on time) + P(Arriving on time | Waking up late) * P(Waking up late)} \\
\\ 
\texttt{Since we don't have the probabilities of waking up on time or late, we can assume that these probabilities are equal (P(Waking up on time) = P(Waking up late) = 0.5) for simplicity. This is a common assumption in causal inference when the probabilities of the treatment (waking up on time or late) are not given.} \\
\\ 
\texttt{Now, we can plug in the values:} \\
\\ 
\texttt{P(Arriving on time) = 0.46 * 0.5 + 0.78 * 0.5 = 0.62} \\
\\ 
\texttt{Now we can use the formula for counterfactuals:} \\
\\ 
\texttt{P(Arriving on time | Waking up late, Counterfactual: Waking up on time) = 0.46 + 0.78 {-} 0.62 = 0.62} \\
\\ 
\texttt{However, this is not the answer we are looking for. We want to know if it would be more likely to see arriving to school on time if Alice had gotten up on time, given that she woke up late. This is equivalent to comparing the counterfactual probability to the actual probability of arriving on time given that Alice woke up late.} \\
\\ 
\texttt{P(Arriving on time | Waking up late) = 0.78} \\
\\ 
\texttt{P(Arriving on time | Waking up late, Counterfactual: Waking up on time) = 0.62} \\
\\ 
\texttt{Since 0.62 < 0.78, it would not be more likely to see arriving to school on time if Alice had gotten up on time, given that she woke up late.} \\
            }
        }
    }
    & \\ \\

    \theutterance \stepcounter{utterance}  
    & & \multicolumn{4}{p{0.6\linewidth}}{
        \cellcolor[rgb]{0.95,0.95,0.95}{
            \makecell[{{p{\linewidth}}}]{
                \texttt{\tiny{[P1$\leftrightarrow$P]}}
                \texttt{{[}SELF{-}FEEDBACK{]}} \\
\texttt{FORMAT: ok} \\
\texttt{RULES: ok} \\
\texttt{STRATEGY: problem} \\
\texttt{REASON: The answer correctly applies the concept of counterfactuals and potential outcomes, but the conclusion drawn from the calculation is incorrect, as the counterfactual probability is actually lower than the actual probability of arriving on time given that Alice woke up late.} \\
            }
        }
    }
    & \\ \\

    \theutterance \stepcounter{utterance}  
    & \multicolumn{4}{p{0.6\linewidth}}{
        \cellcolor[rgb]{0.561,0.737,0.561}{
            \makecell[{{p{\linewidth}}}]{
                \texttt{\tiny{[P1$\rangle$GM]}}
                \texttt{No.} \\
\\ 
\texttt{The counterfactual probability of arriving on time if Alice had gotten up on time, given that she woke up late, is 0.62. This is lower than the actual probability of arriving on time given that Alice woke up late, which is 0.78. Therefore, it would not be more likely to see arriving to school on time if Alice had gotten up on time, given that she woke up late.} \\
            }
        }
    }
    & & \\ \\

    \theutterance \stepcounter{utterance}  
    & & & \multicolumn{2}{p{0.3\linewidth}}{
        \cellcolor[rgb]{0.9,0.9,0.9}{
            \makecell[{{p{\linewidth}}}]{
                \texttt{\tiny{[GM$|$GM]}}
                \texttt{No} \\
            }
        }
    }
    & & \\ \\

    \theutterance \stepcounter{utterance}  
    & & & \multicolumn{2}{p{0.3\linewidth}}{
        \cellcolor[rgb]{0.9,0.9,0.9}{
            \makecell[{{p{\linewidth}}}]{
                \texttt{\tiny{[GM$|$GM]}}
                \texttt{no} \\
            }
        }
    }
    & & \\ \\

    \theutterance \stepcounter{utterance}  
    & & & \multicolumn{2}{p{0.3\linewidth}}{
        \cellcolor[rgb]{0.9,0.9,0.9}{
            \makecell[{{p{\linewidth}}}]{
                \texttt{\tiny{[GM$|$GM]}}
                \texttt{game\_result = WIN} \\
            }
        }
    }
    & & \\ \\

\end{supertabular}
}

\end{document}
